% GENERATED FILE. Edit canonical lessons, then run npm run generate:books.
% canonical-source-hash: fnv1a64:13d3cc40f1541814
% canonical-lessons: MW-C12-hear-mausam, MW-W12-au-matra, MW-C12-mausam, MW-C12-hear-garmi, MW-W12-ga, MW-C12-garmi, MW-R12-weather-first-two, MW-C12-hear-thandi, MW-W12-dda, MW-C12-thandi, MW-R12-weather-new-three, MW-C12-weather-six, MW-R12-script-close

\chapter{Weather, Heat, and Cold: Three New Signs}
\label{ch:mw-weather-six}
\hlchaptermodality{17}

\begin{chapteropening}
\textbf{By the end of this chapter:} \emph{I can add weather, heat, and cold to the earlier wind-cloud-rain map after learning only the au mark, \mw{ग}, and \mw{ड}.}

Three more source-attested weather words begin by ear, add three signs one at a time, and join the earlier map in four separately scored skills without claiming a complete weather exchange.
\end{chapteropening}

\section[mausam]{Hear \emph{mausam} — weather}

\label{lesson:MW-C12-hear-mausam}



\begin{quote}
Recall wind, cloud, and rain in four skills, then write \textbf{\mw{ई}}, \textbf{\mw{कांई}}, and \textbf{\mw{घर}}.
\end{quote}

\subsection*{What to know first}
\begin{quote}
\emph{mausam} — \textbf{weather}
\end{quote}

Listen for the \emph{au} sound. Say the general word once; its new vowel mark stays off the page until the next lesson.

\begin{tcolorbox}[breakable,colback=teal!4,colframe=teal!35!black,title={Before you move on}]
Source: \href{https://www.marwaripathshala.com/marwari-lesson-7-english}{Marwari Pathshala, Lesson 7}.
\end{tcolorbox}
\section[au]{\mw{ौ} — the \emph{au} mark in weather}

\label{lesson:MW-W12-au-matra}



\begin{quote}
Say weather and temple, write \textbf{\mw{म}}, \textbf{\mw{बरसात}}, and \textbf{\mw{कांई}}, then contrast familiar \textbf{\mw{ै}} and \textbf{\mw{ो}}.
\end{quote}

\subsection*{Script — one sign}
\begin{quote}
\textbf{\mw{ौ}} — \emph{au} after a consonant
\end{quote}

\subsection*{Your turn}
Trace \textbf{\mw{ौ}} on \textbf{\mw{म}} once, copy \textbf{\mw{मौ}} once, cover it, and write \textbf{\mw{मौ}} from sound.

\begin{tcolorbox}[breakable,colback=teal!4,colframe=teal!35!black,title={Before you move on}]
Source: \href{https://www.unicode.org/charts/PDF/U0900.pdf}{Unicode Devanagari chart}.
\end{tcolorbox}
\section[mausam]{\mw{मौसम} — weather on the page}

\label{lesson:MW-C12-mausam}



\begin{quote}
Ask and answer the known name exchange, say weather and wind, write \textbf{\mw{मंदिर}}, \textbf{\mw{बादल}}, and \textbf{\mw{थारो}}, then form \textbf{\mw{मौ}} and \textbf{\mw{स}}.
\end{quote}

\subsection*{What to know first}
\begin{quote}
\textbf{\mw{मौ}} + \textbf{\mw{सम}} $\to$ \textbf{\mw{मौसम}} — \emph{mausam} — weather
\end{quote}

\subsection*{Your turn}
Look, cover, wait five seconds, and write \textbf{\mw{मौसम}}. Check the first-syllable vowel mark.

\begin{tcolorbox}[breakable,colback=teal!4,colframe=teal!35!black,title={Before you move on}]
Source: \href{https://www.marwaripathshala.com/marwari-lesson-7-english}{Marwari Pathshala, Lesson 7}.
\end{tcolorbox}
\section[garmī]{Hear \emph{garmī} — heat}

\label{lesson:MW-C12-hear-garmi}



\begin{quote}
Say and write weather, recall rain and the three-place payoff, ask \emph{āp kaiso ho?}, then write the \textbf{\mw{है}} ending from the known name line.
\end{quote}

\subsection*{What to know first}
\begin{quote}
\emph{garmī} — \textbf{heat; hot weather}
\end{quote}

The source also groups this word with summer. Keep the narrow heat meaning for now; the season contrast comes later.

\begin{tcolorbox}[breakable,colback=teal!4,colframe=teal!35!black,title={Before you move on}]
Source: \href{https://www.marwaripathshala.com/marwari-lesson-7-english}{Marwari Pathshala, Lesson 7}.
\end{tcolorbox}
\section[ga]{\mw{ग} — the start of heat}

\label{lesson:MW-W12-ga}



\begin{quote}
Say heat, weather, and hand; write \textbf{\mw{आप} \mw{कैसो} \mw{हो}?}, \textbf{\mw{घ}}, \textbf{\mw{मौ}}, and \textbf{\mw{म्हारो}}.
\end{quote}

\subsection*{Script — one sign}
\begin{quote}
\textbf{\mw{ग}} — \emph{ga}
\end{quote}

Compare \textbf{\mw{ग}} with aspirated \textbf{\mw{घ}}; this word begins with the unaspirated sign.

\subsection*{Your turn}
Trace once, copy once, cover the model, and write \textbf{\mw{ग}} from sound.

\begin{tcolorbox}[breakable,colback=teal!4,colframe=teal!35!black,title={Before you move on}]
Source: \href{https://www.unicode.org/charts/PDF/U0900.pdf}{Unicode Devanagari chart}.
\end{tcolorbox}
\section[garmī]{\mw{गर्मी} — heat on the page}

\label{lesson:MW-C12-garmi}



\begin{quote}
Say heat, weather, wind, and the known fine answer; write \textbf{\mw{हाथ}} and \textbf{\mw{मौसम}}, then form \textbf{\mw{ग}}, \textbf{\mw{र्}}, \textbf{\mw{म}}, and \textbf{\mw{ी}}.
\end{quote}

\subsection*{What to know first}
\begin{quote}
\textbf{\mw{गर्}} + \textbf{\mw{मी}} $\to$ \textbf{\mw{गर्मी}} — \emph{garmī} — heat
\end{quote}

\subsection*{Your turn}
Look, cover, wait five seconds, and write \textbf{\mw{गर्मी}}. Check the virama under \textbf{\mw{र}} and long \textbf{\mw{ी}} after \textbf{\mw{म}}.

\begin{tcolorbox}[breakable,colback=teal!4,colframe=teal!35!black,title={Before you move on}]
Source: \href{https://www.marwaripathshala.com/marwari-lesson-7-english}{Marwari Pathshala, Lesson 7}.
\end{tcolorbox}
\section[Practice]{Retrieve weather and heat}

\label{lesson:MW-R12-weather-first-two}



\begin{quote}
Say money, then write \textbf{\mw{मौ}}, \textbf{\mw{ग}}, \textbf{\mw{बादल}}, \textbf{\mw{है}}, and \textbf{\mw{ू}}.
\end{quote}

\subsection*{Your turn}
Hear the two words in mixed order, give each meaning, read both cards, then write both from sound.

\begin{tcolorbox}[breakable,colback=teal!4,colframe=teal!35!black,title={Before you move on}]
Sources: \href{https://www.marwaripathshala.com/marwari-lesson-7-english}{Marwari Pathshala, Lesson 7}; \href{https://www.unicode.org/charts/PDF/U0900.pdf}{Unicode Devanagari chart}.
\end{tcolorbox}
\section[ṭhaṇḍī]{Hear \emph{ṭhaṇḍī} — cold}

\label{lesson:MW-C12-hear-thandi}



\begin{quote}
Say and write weather and heat, recall cloud, then write rain and money.
\end{quote}

\subsection*{What to know first}
\begin{quote}
\emph{ṭhaṇḍī} — \textbf{cold; cold weather}
\end{quote}

The source also groups this form with winter. Keep the narrow cold meaning for now and contrast it with heat.

\begin{tcolorbox}[breakable,colback=teal!4,colframe=teal!35!black,title={Before you move on}]
Source: \href{https://www.marwaripathshala.com/marwari-lesson-7-english}{Marwari Pathshala, Lesson 7}.
\end{tcolorbox}
\section[ḍa]{\mw{ड} — the middle of cold}

\label{lesson:MW-W12-dda}



\begin{quote}
Say cold, heat, and the known wellbeing answer; write \textbf{\mw{हूं} \mw{ठीक} \mw{हूं}}, then \textbf{\mw{द}}, \textbf{\mw{ठ}}, \textbf{\mw{ग}}, and \textbf{\mw{हवा}}.
\end{quote}

\subsection*{Script — one sign}
\begin{quote}
\textbf{\mw{ड}} — \emph{ḍa}, with the tongue curled back
\end{quote}

Keep \textbf{\mw{ड}} distinct from dental \textbf{\mw{द}}.

\subsection*{Your turn}
Trace once, copy once, cover the model, and write \textbf{\mw{ड}} from sound.

\begin{tcolorbox}[breakable,colback=teal!4,colframe=teal!35!black,title={Before you move on}]
Source: \href{https://www.unicode.org/charts/PDF/U0900.pdf}{Unicode Devanagari chart}.
\end{tcolorbox}
\section[ṭhaṇḍī]{\mw{ठंडी} — cold on the page}

\label{lesson:MW-C12-thandi}



\begin{quote}
Perform the known wellbeing exchange and travel payoff, say cold, heat, and weather, write \textbf{\mw{गर्मी}}, then form \textbf{\mw{ठ}}, \textbf{\mw{ं}}, \textbf{\mw{ड}}, and \textbf{\mw{ी}}.
\end{quote}

\subsection*{What to know first}
\begin{quote}
\textbf{\mw{ठं}} + \textbf{\mw{डी}} $\to$ \textbf{\mw{ठंडी}} — \emph{ṭhaṇḍī} — cold
\end{quote}

\subsection*{Your turn}
Look, cover, wait five seconds, and write \textbf{\mw{ठंडी}}. Check the nasal mark and keep \textbf{\mw{ड}} distinct from \textbf{\mw{द}}.

\begin{tcolorbox}[breakable,colback=teal!4,colframe=teal!35!black,title={Before you move on}]
Source: \href{https://www.marwaripathshala.com/marwari-lesson-7-english}{Marwari Pathshala, Lesson 7}.
\end{tcolorbox}
\section[Practice]{Retrieve weather, heat, and cold}

\label{lesson:MW-R12-weather-new-three}



\begin{quote}
Say the known see-you-later line, then write \textbf{\mw{मौ}}, \textbf{\mw{ग}}, \textbf{\mw{ड}}, and \textbf{\mw{बरसात}}.
\end{quote}

\subsection*{Your turn}
Hear the three words in mixed order, give each meaning, read three cards, then write all three from sound.

\begin{tcolorbox}[breakable,colback=teal!4,colframe=teal!35!black,title={Before you move on}]
Sources: \href{https://www.marwaripathshala.com/marwari-lesson-7-english}{Marwari Pathshala, Lesson 7}; \href{https://www.unicode.org/charts/PDF/U0900.pdf}{Unicode Devanagari chart}.
\end{tcolorbox}
\section[Practice]{Payoff — six weather words in four skills}

\label{lesson:MW-C12-weather-six}



\begin{quote}
Write \textbf{\mw{छ}}, then recall the earlier wind-cloud-rain payoff before adding three words.
\end{quote}

\subsection*{Your turn}
1. Identify six heard words. 2. Produce six words from meaning cues. 3. Match six printed cards to meanings. 4. Write three heard words from each half without a model.

Pass each skill separately.

\begin{tcolorbox}[breakable,colback=teal!4,colframe=teal!35!black,title={Before you move on}]
Source: \href{https://www.marwaripathshala.com/marwari-lesson-7-english}{Marwari Pathshala, Lesson 7}.
\end{tcolorbox}
\section[Practice]{Close the chapter — three new signs}

\label{lesson:MW-R12-script-close}



\begin{quote}
Write \textbf{\mw{ज}}, \textbf{\mw{घ}}, \textbf{\mw{े}}, and \textbf{\mw{हवा}}, then name the four weather-payoff skills.
\end{quote}

\subsection*{Your turn}
From sound alone, write \textbf{\mw{मौ}}, \textbf{\mw{ग}}, \textbf{\mw{ड}}, then \textbf{\mw{मौसम}}, \textbf{\mw{गर्मी}}, and \textbf{\mw{ठंडी}}. Give each word's meaning.

\begin{tcolorbox}[breakable,colback=teal!4,colframe=teal!35!black,title={Before you move on}]
Sources: \href{https://www.unicode.org/charts/PDF/U0900.pdf}{Unicode Devanagari chart}; \href{https://www.marwaripathshala.com/marwari-lesson-7-english}{Marwari Pathshala, Lesson 7}.
\end{tcolorbox}
